\section{Open Questions}
\label{sec:open-questions}

The replication established that the paper's specific claims hold on the
projector-cost / single-target class it analyzes. Five questions remain
genuinely open --- none is answered by the paper, and each has a concrete,
free-endpoint, CPU-tractable probe.

\begin{enumerate}
\item \textbf{Warm-start / meta-learning transfer to structured combinatorial
problems.} The paper proves concentration for the trivial cost
$H_{z}=1-|t\rangle\langle t|$. Whether the concentrated
$(\gamma^{\star},\beta^{\star})$ are a useful warm-start for MaxCut on
random-regular graphs, Sherrington--Kirkpatrick, etc.\ is the natural
operational question. \emph{Next steps:} run QAOA at $p=1..5$ on 3-regular
MaxCut, $n=8..20$, comparing cold random seeds vs the paper's
$\bigl(\pi(n+2)/(n+4),\pi/(n+2)\bigr)$ vs an edge-weight-rescaled variant.
Report median iterations-to-convergence and final approximation ratio.
qiskit-aer statevector, CPU, free.

\item \textbf{When does concentration fail?} No published characterization
of the failure boundary exists. Does concentration break for heterogeneous
target ensembles / full-spectrum $H_{z}$? Can breakdown be predicted from
one scalar (spectral gap $\Delta$, effective rank $r$)? \emph{Next steps:}
sweep synthetic $H_{z}=\sum_{k}\lambda_{k}|k\rangle\langle k|$ with
controlled $\Delta$ and $r$; run QAOA at $p=1..3$; measure $\Delta^{2}$
vs $n$; fit the crossover where $\Delta^{2}$ saturates. Statevector,
$n\le14$, free.

\item \textbf{Symmetry-informed initialization $\otimes$ concentration.}
Does composing concentration with problem-symmetry seeds (Shaydulin~2021,
Egger~2021 warm-start) stack multiplicatively, additively, or destructively?
\emph{Next steps:} on $\mathbb{Z}_{2}$-symmetric MaxCut instances (bipartite,
regular), seed with concentrated angles plus an automorphism-group rotation;
measure iterations-to-99\%-optimum vs each seed alone. $p=1..3$, $n=8..16$,
statevector, free.

\item \textbf{Noise-robust parameter transfer.} The paper is exact statevector.
Whether concentrated angles remain near-optimal under depolarizing / $T_{1}$
/ $T_{2}$ noise, or whether every hardware calibration needs its own
optimization, determines NISQ usefulness. \emph{Next steps:} sweep depolarizing
$p_{\text{err}} \in \{0, 10^{-4}, 10^{-3}, 10^{-2}\}$ with qiskit-aer
\texttt{AerSimulator}; re-optimize at each; plot
$|\theta_{\text{noisy}} - \theta_{\text{conc}}|$ vs $p_{\text{err}}$.
$n=6..10$, $p=1..3$, free.

\item \textbf{ML-boosted concentration prediction.} If the target manifold
is low-dimensional in $(n,p,\text{problem-features})$, a small MLP / KRR /
GP should predict $(\gamma^{\star},\beta^{\star})$ directly, replacing
optimization. \emph{Next steps:} build 300 solved
$(n,p,\text{edge-density},\text{spectral-gap})\!\to\!(\gamma^{\star},\beta^{\star})$
tuples for MaxCut $n=6..14$, $p=1..3$; train sklearn \texttt{KernelRidge}
with RBF kernel; evaluate leave-one-instance-class-out MAE and downstream
approximation-ratio loss. CPU, free.
\end{enumerate}
